\subsection{Configuración de Experimentos}

\begin{table}[H]
  \centering
  \caption{Diseño experimental completo}
  \label{tab:diseno}
  \begin{tabular}{ll}
    \toprule
    \textbf{Factor} & \textbf{Niveles} \\
    \midrule
    Instancias       & \texttt{small} (100), \texttt{medium} (1000), \texttt{large} (10000) \\
    Tamaños de pob.  & 1024, 4096, 16384 \\
    Variantes        & \texttt{sequential}, \texttt{cuda\_basic}, \texttt{cuda\_optimized} \\
    Generaciones     & 300 (fijo) \\
    Tamaño de bloque & 128 (EXP 1, 3, 4); 32, 64, 128, 256 (EXP 2)\\
    Semillas         & 42, 43, 44, 45, 46, 47, 48, 49, 50, 51 (todos los EXP) \\
    \bottomrule
  \end{tabular}
\end{table}

\subsection{Descripción de los Cuatro Experimentos}

\paragraph{EXP 1 — Diseño principal (270 ejecuciones).}
Compara las tres variantes en los tres tamaños de instancia y los tres tamaños de
población, con 10 semillas por celda. Mide tiempo total, tiempo de kernels,
porcentaje de soluciones factibles y calidad de solución.

\paragraph{EXP 2 — Efecto del tamaño de bloque (80 ejecuciones).}
Fija instancia \texttt{large} y población 4096; varía el bloque
(\{32, 64, 128, 256\}) para las dos variantes CUDA con 10 semillas. Permite
identificar la configuración que mejor aprovecha la ocupancia de los SMs.

\paragraph{EXP 3 — Análisis de \emph{speed-up} (270 ejecuciones).}
Para cada combinación instancia × población × semilla de EXP 1
mide el tiempo de las tres variantes y calcula \emph{speed-up} =
$T_{\text{seq}} / T_{\text{GPU}}$.

\paragraph{EXP 4 — Efecto del tamaño de población (90 ejecuciones).}
Fija instancia \texttt{large}; varía población y variante con 10 semillas. Relaciona
calidad de solución y tiempo con el número de individuos.

\subsection{Reproducibilidad}

Todas las ejecuciones usan semillas explícitas tanto en la CPU
(\texttt{std::mt19937(seed)}) como en la GPU (\texttt{curand\_init(seed, ind, 0,
\&states[ind])}). Los resultados se almacenan en archivos CSV en el directorio
\texttt{results/} con el nombre de cada experimento.
